\subsection{Dostupnost modelu}
\label{subsec:criteria-availability}

Dostupností se v tomto kontextu rozumí několik vzájemně propojených otázek:

\begin{enumerate}
    \item \textbf{Způsob přístupu} -- je model dostupný jako cloudové API nebo jako open-source model ke stažení a lokálnímu nasazení?
    Cloudové API nabízí snadnou integraci bez hardwarových nároků, ale zavádí závislost na externím poskytovateli.
    Open-source modely umožňují lokální nasazení a plnou kontrolu nad zpracováním dat.

    \item \textbf{Licenční podmínky} -- pro školní použití v rámci instituce je nutné ověřit, zda licence modelu takové použití umožňuje.
    Mnoho modelů je dostupných pod striktními licencemi (Apache~2.0, MIT), ale některé mají omezení pro komerční použití nebo vyžadují uzavření samostatné licenční smlouvy.

    \item \textbf{Geografická dostupnost a datová centra} -- jak bylo zmíněno v \secref[Sekci]{subsubsec:gdpr-data-transfer} o ochraně osobních údajů, je důležité, aby poskytovatel modelu měl datová centra v regionech, které jsou v souladu s GDPR, v případě možného úniku dat.

    \item \textbf{Stabilita a podpora} -- API musí být spolehlivé a třetí strana musí poskytovat průběžnou podporu.
    Pro školní použití je důležitá i předvídatelnost cenového modelu a garancí dostupnosti (SLA).
\end{enumerate}

\endinput